% sec3b.tex — c07｜原文 p23–25｜§3.3 class S 对偶与 TQFT（续）、Schur 极限、推广到更高秩
% 编号公式 (3.20)–(3.35) 共 16 条；无编号环境、无表格。
% [原文 p.23]
在退化极限下，四穿刺球面退化并分裂成三个三穿刺球面。根据 Gaiotto，一个三穿刺球面是相对于三个穿刺点处 $\SU(2)$ 的三重基础表示中的半超多重态。带三个穿刺点的球面没有任何模，这与自由半超多重态理论没有任何耦合相符。可以取三种不同的退化极限，即分别把 $a$ 与 $b$、$a$ 与 $c$、或 $a$ 与 $d$ 靠近。在每个退化极限下我们都得到一个弱耦合理论，但其中半超多重态分别处于 $\SO(8)$ 整体对称性的 $\mathbf{8}_v$、$\mathbf{8}_s$ 与 $\mathbf{8}_c$ 表示。只看 $\SU(2)^4$ 子群，该对偶可以看作一种交叉对称性。

更一般地，由于指标与规范耦合无关，因而与 Riemann 曲面的复结构无关，它必须由某个拓扑场论来计算。抽象地说，一个拓扑场论可以通过给出三点函数与两点函数（即传播子）来指定。

让我们回顾这是如何运作的。我们把三穿刺球面的指标参数化为 $\sI(a,b,c)$，其中 $a,b,c$ 是三个 $\SU(2)$ 对称性 Cartan 的逸度。另一方面，我们可以容易地写下与圆柱相对应的「传播子」，
\[
\eta(a,b)=\Delta(a)\sI_V(a)\,\delta(a,b^{-1}) ,
\tag{3.20}
\]
其中 $\Delta(a)$ 是 Haar 测度，$\sI_V(a)$ 是矢量多重态的指标，后者已知显式形式。一般 class S 理论的指标可以用这些基本组元的指标表示。最简单的例子是：用一个圆柱把两个三穿刺球面粘起来，就得到四穿刺球面的指标，
\[
\sI(a,b,c,d)=\oint\frac{\dd z}{2\pi\ii}\oint\frac{\dd x}{2\pi\ii}\sI_H(a,b,z)\,\eta(z,x)\,\sI_H(z,c,d)
\tag{3.21}
\]
\[
=\oint\dd z\,\Delta(z)\,\sI_H(a,b,z)\,\sI_V(z)\,\sI_H(z,c,d) .
\tag{3.22}
\]
若把指标按方便的函数基 $\{f^\alpha(a)\}$ 展开，该基以 $\SU(2)$ 表示 $\{\alpha\}$ 标号\footnote{对 type A 理论，$\{f^\alpha(a)\}$ 是其自变量的对称函数，这些自变量是与 $\SU(k)$ 的 Cartan 生成元对偶的逸度。更一般地，对 type D 与 E 理论，$\{f^\alpha(a)\}$ 在相应的 Weyl 群下不变。}，我们就可以给每个三穿刺球面指定「结构常数」$C_{\alpha\beta\gamma}$，给每个传播子指定度规 $\eta^{\alpha\beta}$，
\[
\sI(a,b,c)=\sum_{\alpha,\beta,\gamma}C_{\alpha\beta\gamma}\,f^\alpha(a)\,f^\beta(b)\,f^\gamma(c)
\tag{3.23}
\]
\[
\eta^{\alpha\beta}=\oint\frac{\dd a}{2\pi\ii}\oint\frac{\dd b}{2\pi\ii}\eta(a,b)\,f^\alpha(a)\,f^\beta(b) .
\tag{3.24}
\]
指标在曲面的不同分解方式下不变，等价于
% [原文 p.24]
说 $C_{\alpha\beta\gamma}$ 与 $\eta^{\alpha\beta}$ 定义了一个二维拓扑场论\footnote{我们在这个词的用法上有些不严格。按 Atiyah 的公理化定义，TQFT 应具有有限维状态空间，而在我们的情形状态空间是无穷维的。具有无穷维状态空间的二维拓扑理论中理解得最清楚的例子是二维 Yang–Mills 理论的零面积极限 [35,36]（综述见 [37]）。幸运的是，与指标相伴的二维拓扑理论与二维 Yang–Mills 密切相关。}。关键性质是结合性，
\[
C_{\alpha\beta\zeta}\,C^{\zeta}_{\ \gamma\delta}=C_{\alpha\gamma\zeta}\,C^{\zeta}_{\ \beta\delta} ,
\tag{3.25}
\]
其中指标用度规 $\eta^{\alpha\beta}$ 上升、用逆度规 $\eta_{\alpha\beta}$ 下降。

很自然的选择是让完备函数集 $\{f^\alpha(a)\}$ 在传播子中出现的测度下正交归一，
\[
\oint\frac{\dd a}{2\pi\ii}\Delta(a)\sI_V(a)\,f^\alpha(a)f^\beta(a)=\delta^{\alpha\beta} .
\tag{3.26}
\]
于是度规 $\eta^{\alpha\beta}$ 是平庸的，
\[
\eta^{\alpha\beta}=\delta^{\alpha\beta} .
\tag{3.27}
\]
条件 (3.26) 仍留下相当大的自由度，因为无穷多个彼此由正交变换联系起来的函数基都满足它。真正的简化出现在我们能找到一个显式基 $\{f^\alpha(a)\}$，使得结构常数是对角的，
\[
C_{\alpha\beta\gamma}\neq0\quad\to\quad\alpha=\beta=\gamma .
\tag{3.28}
\]
此时结合性 (3.25) 自动成立。对满足 (3.25) 的结构常数，总可以找到使它们对角的基。困难在于以具体形式描述该基。

一般而言，传播子中出现的测度很复杂，没有现成的显式正交归一函数集。我们觉得考虑如下拟设（ansatz）非常有用
\[
f^\alpha(a)=\mathcal{K}(a)\,P^\alpha(a) ,
\tag{3.29}
\]
其中 $\mathcal{K}(a)$ 是某个函数。显然，由 (3.26)，函数 $\{P^\alpha(a)\}$ 在新测度 $\hat\Delta(a)$ 下正交归一，
\[
\oint\frac{\dd a}{2\pi\ii}\hat\Delta(a)\,P^\alpha(a)P^\beta(a)=\delta^{\alpha\beta},\qquad
\hat\Delta(a)\equiv\sI_V(a)\mathcal{K}(a)^2\Delta(a) .
\tag{3.30}
\]
回忆 $\Delta(a)$ 是 Haar 测度。问题的关键在于巧妙地选取 $\mathcal{K}(a)$，使得 $\hat\Delta(a)$ 是简单的已知测度，且正交归一基 $\{P^\alpha(a)\}$ 是满足 (3.28) 的显式函数集。

一旦已知对角基 $\{f^\alpha(a)\}$ 与结构常数 $C_{\alpha\alpha\alpha}$，就容易计算与亏格 $g$、带 $s$ 个穿刺点的曲面相伴的 SCFT 的指标。这样的曲面可以通过粘合 $2g-2+s$ 个三穿刺球面得到，于是
\[
\sI_{g,s}(a_1,a_2,\ldots,a_s)=\sum_\alpha\big(C_{\alpha\alpha\alpha}\big)^{2g-2+s}\prod_{I=1}^{s}f^\alpha(a_I) .
\tag{3.31}
\]
\paragraph{Schur 极限} 在 Schur 极限下最容易辨认正交归一函数基。在此极限下，
\[
\sI_V(a)=PE\Big[\frac{-2q}{1-q}(a^2+1+a^{-2})\Big],\qquad
\sI_H(a,b,c)=PE\Big[\frac{\sqrt q}{1-q}\Big(a+\frac1a\Big)\Big(b+\frac1b\Big)\Big(c+\frac1c\Big)\Big] .
\tag{3.32}
\]
% [原文 p.25]
若取 $\mathcal{K}(a)=1/\sqrt{\sI_V(a)}$，则 $\hat\Delta(a)=\Delta(a)$。关于标准 Haar 测度正交归一的函数基是特征标，也称为 Schur 多项式。确实，三点函数在该基下是对角的 [38]，即
\[
\sI_H(a,b,c)=\frac{C(q)}{\sqrt{\sI_V(a)\sI_V(b)\sI_V(c)}}\sum_\alpha\frac{1}{\dim_q\alpha}\chi_\alpha(a)\chi_\alpha(b)\chi_\alpha(c) .
\tag{3.33}
\]
由于与该指标极限相伴的正交多项式是 Schur 多项式，该极限称为 Schur 极限。

在 Hall–Littlewood 与 Macdonald 极限中，出现的正交多项式分别是 Hall–Littlewood 与 Macdonald 多项式。对完整指标，$P_\alpha(a)$ 不是多项式，而是与 Ruijsenaars–Schnider 可积模型相关的复杂函数。该模型位于 Calogero–Moser 层级的顶端。

\paragraph{推广到更高秩} 如前所述，$\sN=2$ 的 $\SU(2)$ 规范理论由六维秩 1 的 $(2,0)$ SCFT 紧化得到。若把更高秩的 $(2,0)$ 理论紧化在 Riemann 曲面上，会得到类似的网络，但与秩 1 情形不同，这些理论大多没有拉氏量描述。特别地，对秩 $N$，与三穿刺球面相伴的理论称为 $T_N$ 理论。$T_2$ 理论只是自由超多重态的理论，而 $T_{N>2}$ 是非平庸的强耦合 SCFT。这使人无法走计算指标的直接路线。然而，由于拓扑三点函数与两点函数的表达式具有可自然提升到更高秩的形式，人们就得到了整个强耦合理论族的超共形指标的自然猜想。同样，在 Schur 极限下举例最简单。由于 $T_2$ 理论的指标为
\[
\sI_{T_2}(a,b,c)=\frac{C_2(q)}{\sqrt{\sI_V(a)\sI_V(b)\sI_V(c)}}\sum_\alpha\frac{1}{\dim_q\alpha}\chi_\alpha(a)\chi_\alpha(b)\chi_\alpha(c) ,
\tag{3.34}
\]
这一形式到一般 $N$ 的提升将是
\[
\sI_{T_N}(a_i,b_i,c_i)=\frac{C_N(q)}{\sqrt{\sI_V(a_i)\sI_V(b_i)\sI_V(c_i)}}\sum_R\frac{1}{\dim_q R}\chi_R(a_i)\chi_R(b_i)\chi_R(c_i) ,
\tag{3.35}
\]
其中 $R$ 遍历 $\SU(N)$ 的所有表示，$a_i,b_i,c_i$ 是与整体对称性 $\SU(N)_i$ 相伴的逸度。

对 Macdonald 极限 [38] 乃至完整指标 [39] 也已作出类似的猜想。

% [Tutor Note: 物理：class S 理论即六维 $(2,0)$ 理论紧化在 Riemann 曲面上；三穿刺球面是构建一切的基本积木。]
% [Tutor Note: 易错：(3.20) 传播子中是 $\delta(a,b^{-1})$ 而不是 $\delta(a,b)$；$\SU(2)$ 特征标以 $a\leftrightarrow a^{-1}$ 配对。]
% [Tutor Note: 动机：指标与耦合无关，故只依赖曲面的复结构，因此必是拓扑场论（TQFT）的量——这是本节的中心思想。]
% [Tutor Note: 易错：(3.21) 与 (3.22) 各占一个编号：前者保留 $\eta$，后者把 $\eta(z,x)$ 代入并使用 $\delta(z,x^{-1})$ 消去一个积分。]
% [Tutor Note: 术语：结构常数 $C_{\alpha\beta\gamma}$、度规 $\eta^{\alpha\beta}$、结合性 (3.25)，正是 Atiyah 二维 TQFT 公理的具体化。]
% [Tutor Note: 易错：(3.26) 的正交归一测度是 $\Delta(a)\sI_V(a)$，不是单独的 Haar 测度；后者只在 Schur 极限的 $\hat\Delta=\Delta$ 特例中出现。]
% [Tutor Note: 易错：(3.24) 定义 $\eta^{\alpha\beta}$（指标在上），(3.25) 用 $\eta^{\alpha\beta}$ 上升、$\eta_{\alpha\beta}$ 下降；原文把二者称作 metric 与 inverse metric。]
% [Tutor Note: 物理：对角化结构常数 (3.28) 使结合性自动满足，问题化为「找基」；Schur 极限下这组基就是 Schur 多项式。]
% [Tutor Note: 易错：(3.29) 的 $\mathcal{K}(a)$ 是待定函数，不是指标；它的作用是把复杂测度吸收进 $\hat\Delta(a)$。]
% [Tutor Note: 物理：(3.31) 的幂次 $2g-2+s$ 正是粘合三穿刺球面所需的个数，由此得到任意亏格曲面的指标。]
% [Tutor Note: 易错：(3.32) 中 $PE$ 的参数分别是 $(-2q/(1-q))$ 与 $(\sqrt q/(1-q))$，指数里 $a$ 的组合不能漏掉 $a^{-2}$ 与 $a^{-1}$。]
% [Tutor Note: 文献：三点函数对角性见 [38]；$T_N$ 指标猜想见 [38][39]；与二维 Yang–Mills 的联系见 [35–37]。]
% [Tutor Note: 物理：$T_{N>2}$ 没有拉氏量，其指标只能靠 TQFT 结构「提升」来猜——这是 class S 指标计算威力的体现。]
